\documentclass[hyperref=unicode, presentation,10pt]{beamer}

\usepackage[absolute,overlay]{textpos}
\usepackage{array}
\usepackage{graphicx}
\usepackage{adjustbox}
\usepackage{mhchem}
\usepackage{chemfig}
\usepackage{caption}

%dělení slov
\usepackage{ragged2e}
\let\raggedright=\RaggedRight
%konec dělení slov

\addtobeamertemplate{frametitle}{
	\let\insertframetitle\insertsectionhead}{}
\addtobeamertemplate{frametitle}{
	\let\insertframesubtitle\insertsubsectionhead}{}

\makeatletter
\CheckCommand*\beamer@checkframetitle{\@ifnextchar\bgroup\beamer@inlineframetitle{}}
\renewcommand*\beamer@checkframetitle{\global\let\beamer@frametitle\relax\@ifnextchar\bgroup\beamer@inlineframetitle{}}
\makeatother
\setbeamercolor{section in toc}{fg=red}
\setbeamertemplate{section in toc shaded}[default][100]

\usepackage{fontspec}
\usepackage{unicode-math}

\usepackage{polyglossia}
\setdefaultlanguage{czech}

\def\uv#1{„#1“}

\mode<presentation>{\usetheme{default}}
\usecolortheme{crane}

\setbeamertemplate{footline}[frame number]

\title[Crisis]
{C2062 -- Anorganická chemie II}

\subtitle{Demonstrační pokusy -- p-blok}
\author{Zdeněk Moravec, hugo@chemi.muni.cz \\ \adjincludegraphics[height=60mm]{../IUPAC_PSP.jpg}}
\date{}

\begin{document}
\begin{frame}
	\titlepage
\end{frame}

\section{Demonstrační pokusy -- p-blok}

\frame{
	\frametitle{}
	\vfill
	\textbf{Redukční účinky cínatých iontů}
	\begin{itemize}
		\item \ce{2 HgCl2 + Sn^{2+} -> Hg2Cl2 + Sn^{4+} + 2 Cl^-} (bílá)
		\item \ce{Hg2Cl2 + Sn^{2+} -> Hg^0 + Sn^{4+} + 2 Cl^-} (šedá)
	\end{itemize}
	\vspace{5mm}
	\hrule
	\vspace{5mm}

	\textbf{Analytické reakce olovnatých iontů}
	\begin{itemize}
		\item \ce{Pb^{2+} + 2 Cl^- -> PbCl2} (bílá)
		\item \ce{PbCl2 + 2 Cl^- -> [PbCl4]^{2-}}
		\item \ce{[PbCl4]^{2-} + 2 Cl^- -> [PbCl6]^{4-}}
		\item \ce{Pb^{2+} + 2 SO$_4^{2-}$ -> PbSO4}
		\item \ce{Pb^{2+} + H2SO4 -> Pb(HSO4)2}
		\item \ce{Pb^{2+} + 2 OH^- -> Pb(OH)2}
		\item \ce{Pb(OH)2 + 2 OH^- -> [Pb(OH)4]^{2-}}
	\end{itemize}
	\vfill
}

\frame{
	\frametitle{}
	\vfill
	\textbf{Chromová žluť}
	\begin{itemize}
		\item Olovo je součástí několika pigmentů, vzhledem k toxicitě olova se ale už dnes nepoužívají.
		\item Chromová žluť, \ce{PbCrO4}, v přírodě se nachází jako minerál \textit{krokoit}.\footnote[frame]{\href{https://mineraly.sci.muni.cz/sulfaty/krokoit.html}{Krokoit}} Má silné oxidační účinky.\footnote[frame]{\href{https://pubchem.ncbi.nlm.nih.gov/compound/24460}{Lead chromate -- PubChem}}
		\item \ce{Pb^{2+} + CrO$_4^{2-}$ -> PbCrO4}
	\end{itemize}

	\begin{columns}
		\begin{column}{.5\textwidth}
			\begin{figure}
				\adjincludegraphics[width=.7\textwidth]{img/Lead_chromate.jpg}
				\caption*{Chromová žluť.\footnote[frame]{Zdroj: \href{https://commons.wikimedia.org/wiki/File:Lead_chromate.JPG}{FK1954/Commons}}}
			\end{figure}
		\end{column}
		\begin{column}{.5\textwidth}
			\begin{figure}
				\adjincludegraphics[width=.85\textwidth]{img/HCS_bus49.jpg}
				\caption*{Žlutý školní autobus.\footnote[frame]{Zdroj: \href{https://commons.wikimedia.org/wiki/File:HCS_bus49.JPG}{William Grimes/Commons}}}
			\end{figure}
		\end{column}
	\end{columns}
	\vfill
}

\frame{
	\frametitle{}
	\vfill
	\begin{columns}
		\begin{column}{0.6\textwidth}
			\begin{itemize}
				\item \textbf{Jodid olovnatý}, \ce{PbI2}, je žlutá pevná látka.
				\item Bromid i jodid můžeme připravit srážením olovnatých solí, u jodidu jde o oblíbený pokus - \textit{Zlatý déšť}.\footnote[frame]{\href{https://cs.wikibooks.org/wiki/Chemick\%C3\%A9_pokusy/Zlat\%C3\%BD\%20d\%C3\%A9\%C5\%A1\%C5\%A5}{Zlatý déšť}}
				\item \ce{Pb(NO3)2 + 2 KI -> PbI2 + 2 KNO3}
				\item \ce{PbI2} je prekurzorem pro výrobu solárních článků založených na perovskitu \ce{CH3NH3PbI3}, které mají vysokou účinnost.
			\end{itemize}
		\end{column}
		\begin{column}{0.4\textwidth}
			\begin{figure}
				\adjincludegraphics[width=.8\textwidth]{img/Iodide_lead(II).jpg}
				\caption*{Jodid olovnatý.\footnote[frame]{Zdroj: \href{https://commons.wikimedia.org/wiki/File:Iodide_lead(II).jpg}{Stefano sct/Commons}}}
			\end{figure}
		\end{column}
	\end{columns}
	\vfill
}

\frame{
	\frametitle{}
	\vfill
	\textbf{Pyroforické olovo}
	\begin{itemize}
		\item Olovo je v jemně práškovém stavu samozápalné.\footnote[frame]{\href{https://www.ped.muni.cz/wchem/sm/hc/anorglab/soubory/navody/32.htm}{Příprava pyroforického olova}}
		\item Jemný prášek můžeme získat termickým rozkladem vinanu olovnatého:
		\item \ce{Pb(C4H2O6) -> Pb + CO2 + 3 CO + H2O}
		\item Práškové olovo pak prudce reaguje se vzdušným kyslíkem za vzniku oxidu olovnatého (příp. olovnato-olovičitého, \ce{Pb3O4})
		\item \ce{2 Pb + O2 -> 2 PbO}
	\end{itemize}

	\begin{figure}
		\adjincludegraphics[height=.35\textheight]{img/Lead_tartrate.png}
		\caption*{Vinan olovnatý}
	\end{figure}
	\vfill
}

\frame{
	\frametitle{}
	\vfill
	\textbf{Saturnův strom}
	\begin{columns}
		\begin{column}{.7\textwidth}
			\begin{itemize}
				\item Protože má zinek zápornější elektrodový potenciál než olovo, je schopen jej vytěsnit ze sloučenin.
				\item \ce{Pb(NO3)2 + Zn -> Pb + Zn(NO3)2}
			\end{itemize}
		\end{column}
		\begin{column}{.3\textwidth}
			\begin{tabular}{|l|r@{,}l|}
				\hline
				\textbf{Elektroda} & \multicolumn{2}{c|}{\textbf{E$^0$ [V]}} \\\hline
				Pb/Pb$^{2+}$ & -0 & 126 \\\hline
				Zn/Zn$^{2+}$ & -0 & 763 \\\hline
			\end{tabular}
		\end{column}
	\end{columns}

	\vspace{5mm}
	\hrule
	\vspace{5mm}

	\textbf{Oxidační účinky \ce{PbO2}}
	\begin{itemize}
		\item Oxid olovičitý, \ce{PbO2}, je velmi silné oxidační činidlo.
		\item Dokáže oxidovat manganaté soli až na manganistan:
		\item \ce{5 PbO2 + 2 Mn^{2+} + 4 H^+ -> 5 Pb^{2+} + 2 MnO$_4^-$ + 2 H2O}
	\end{itemize}
	\vfill
}

\frame{
	\frametitle{}
	\vfill
	\begin{itemize}
		\item Arsen je možno dokázat i \textit{Gutzeitovým testem}.\footnote[frame]{\href{https://doi.org/10.1039/AN9012600181}{The Gutzeit test for arsenic}}
		\item Jeho autorem je německý chemik Max Adolf Gutzeit (1847--1915).
		\item Arsan připravíme podobně jako u Marshova testu:
		\item \ce{6 Zn + 12 HCl + As2O3 -> 2 AsH3 + 6 ZnCl2 + 3 H2O}
		\item Ten je veden na filtrační papír navlhčený koncentrovaným roztokem \ce{AgNO3}:
		\item \ce{AsH3 + 6 AgNO3 -> Ag3As.3AgNO3 + 3 HNO3}
		\item Vzniklý arsenid následně hydrolyzuje na kovové stříbro:
		\item \ce{Ag3As.3AgNO3 + 3 H2O -> 6 Ag + H3AsO3 + 3 HNO3}
	\end{itemize}
	\vfill
	}

\frame{
	\frametitle{}
	\vfill
	\textbf{Woodův kov}
	\begin{itemize}
		\item Woodův kov je nízkotající slitina cínu, olova, bismutu a kadmia. Její teplota tání je v závislosti na složení 60 až 70 °C, proto ji můžeme roztavit s využitím vodní lázně.\footnote[frame]{\href{https://cs.wikibooks.org/wiki/Chemické_pokusy/Tání_Woodova_kovu}{Tání Woodova kovu}}
	\end{itemize}

	\vspace{5mm}
	\hrule
	\vspace{5mm}

	\textbf{Příprava červeného selenu}
	\begin{itemize}
		\item Selen můžeme získat redukcí seleničitanu pomocí oxidu siřičitého:
		\item \ce{8 SeO$_3^{2-}$ + 16 SO2 + 8 H2O -> Se8 + 16 SO$_4^{2-}$ + 16 H^+}
	\end{itemize}
	\vfill
}

\end{document}